\section{Introduction}
\label{sec:intro}

Earnings conference calls are one of the few recurring events at which management
speaks at length, semi-spontaneously, about a public company---in both \emph{what}
is said (the transcript) and \emph{how} it is said (the audio). A line of research
beginning with \citet{qin2019what} and continuing through HTML
\citep{yang2020html}, MAEC \citep{li2020maec}, KeFVP \citep{niu2023kefvp}, and ECC
Analyzer \citep{cao2024ecc} claims that multimodal models over this data predict
post-call stock volatility better than classical methods.

Two recent observations cast doubt on how much of that progress is real. First,
\citet{scss2024} (``Same Company, Same Signal'') show that transcript
representations used across this literature predominantly encode \emph{company
identity} rather than content: same-ticker transcripts cluster tightly, and
training-free baselines built on a company's own past volatility beat all
transcript-based models on their evaluation. Because volatility is strongly
autocorrelated and company-specific, a model that learns ``which company is this''
inherits most of the apparent skill. Second, a 2025 survey of financial NLP
\citep{finnlp2025survey} finds that only $\sim$14\% of reported results reproduce
exactly, attributing the gap to dead data links, API drift, and contamination.

This paper's premise is that both problems are properties of the \emph{evaluation
substrate} the field shares, and that the substrate can be repaired. The lineage
above evaluates on a small number of label and split files that are passed from
paper to paper essentially unchanged. Those files (\S\ref{sec:substrate}) impose no
embargo between training and test windows even though targets run 30 days past
each call; they place most test calls' companies in the training set, and no paper
reports the ticker-disjoint condition; they carry degenerate label cells; and they
encode a horizon convention (calendar days) that the downstream literature
describes as trading days. On such a substrate a model cannot be told apart from a
memory of which company it is looking at. Meanwhile the corpus with the best
license and the only raw audio, FinCall-Surprise \citep{fincall2025}, ships with
\emph{no} ticker, company, date, or speaker metadata at all---so the identity
controls that would expose the confound cannot even be run on it without first
reconstructing the identity the dataset omits.

We therefore build and release \ecvol-bench, an open benchmark that supplies what
the substrate lacks, and we use it to re-examine the task's central claims. The
benchmark is the contribution; the empirical results are the evidence for why it
is needed. Concretely, we ask four research questions:
\begin{description}
  \item[RQ1] Does earnings-call \emph{content} (text and/or audio) add predictive
    signal for post-call volatility beyond past-volatility persistence (HAR-RV,
    GARCH)?
  \item[RQ2] Is multimodal (audio $+$ text) better than the best unimodal model,
    after controlling for past volatility?
  \item[RQ3] Do the published models of this lineage, re-run on their own data,
    retain their reported advantage under embargoed, \tickerdisjoint{} evaluation?
  \item[RQ4] Do any positive results survive (a) ticker-identity controls,
    (b)~\tickerdisjoint{} splits, and (c) evaluation on post-LLM-knowledge-cutoff
    calls?
\end{description}

The study was pre-registered so that no single outcome strands the work; the
modelling stages, controls, and confirmatory tests were fixed before any model
result existed and are recorded in a public decision log. RQ1--RQ4 are all
answered here; every number in the paper regenerates from a manifested run.

\paragraph{Contributions.}
\begin{enumerate}
  \item \textbf{\ecvol-bench}, an open multimodal volatility benchmark over
    FinCall-Surprise (primary) and MAEC (secondary): reconstructed identity and
    dates for FinCall-Surprise (92.9\% of calls; 95.3\% of earnings calls; 60/60
    manual audit), three target families (level $v$, change \dv, HAR-residual)
    under both the calendar-day and trading-day horizon conventions, temporal and
    \tickerdisjoint{} splits with a 30-trading-day embargo, SHA-256 manifests, and
    a reproducible harness---released in full (\S\ref{sec:benchmark}).
  \item \textbf{A substrate audit} of the label and split files the prior
    literature shares: embargo, ticker overlap, label defects, and the horizon
    convention, regenerable from one command (\S\ref{sec:substrate},
    Table~\ref{tab:substrate}).
  \item \textbf{Honest baselines and an identity-control suite}---persistence,
    EWMA, HAR-RV, GARCH(1,1), and ticker-fixed-effect models for every horizon,
    split, and target; identity probes; transcript- and audio-shuffle controls;
    \tickerdisjoint{} evaluation; a gender-confound analysis---and what they show
    when modern open-weight text and audio representations are run through them on
    consumer hardware (\S\ref{sec:results}).
  \item \textbf{Measured call timing}: the earnings-release timestamp of every
    call from EDGAR 8-K filings (89\% of FinCall, 69\% of MAEC, 99.8\% of
    Earnings25), showing that the call-date anchor used by every released label set
    places the reaction session in the pre-window for most calls, with a
    sensitivity study (\S\ref{sec:timing}).
  \item \textbf{A reproduction study} of the runnable models in this lineage on
    their own data under the new controls, with a code-availability audit of the
    seventeen papers (one runnable as released, seven partial, nine none; \S\ref{sec:results-repro},
    Table~\ref{tab:repro}, Table~\ref{tab:code-audit}).
  \item \textbf{Post-cutoff evaluation}: the audio ladder re-run on 2025 recordings
    stratified by bitrate, and the frozen pipeline evaluated on 475 calls held after
    the knowledge cutoff of every model used (\S\ref{sec:results-clean-audio},
    \S\ref{sec:results-lookahead}).
\end{enumerate}
